\section{Component contrasts}
\label{sec:components}

The information-channel ablations did not establish stable incremental effects. Memory minus scalar had a median difference of 0.000290 and a mean of 0.001703. Twelve units were positive, four were tied, and eight were negative. Three task medians were positive, and the task-cluster interval of [$-$0.001174, 0.004862] crossed zero. The exact task sign-flip p-value was 0.625, with Holm-adjusted $p=1.0$.

Diagnostic minus memory was zero at the median and 0.000355 at the mean. Seventeen of the 24 unit differences were tied, four were positive, and three were negative. Only T12 had a positive task median. The task-cluster interval was [0.000000, 0.000000], consistent with the preponderance of zero unit effects in the resampled 24-unit medians. The exact task sign-flip and Holm-adjusted p-values were both 1.0.

These outcomes leave the independent value of cumulative memory and compact diagnostics unresolved within the present protocol. The positive end-to-end comparisons against random and fixed references show that the complete memory-equipped process can outperform those references while the incremental contribution over scalar feedback remains difficult to separate. Tied champions and the six-task test resolution limit what can be inferred from the component nulls; they do not demonstrate that either information channel has zero effect under every task, budget, or harness.
